% methods.tex (국문 v14, 2026-10-20): 자료와 방법(결과 앞). 번호 없는 소절, 수식 (1)-(13).
\section{자료와 방법}\label{sec:methods}

\subsection{연구 지역과 활동층 관측}
여러 해에 걸친 융해 깊이 관측이 조밀한 세 지역, 알래스카 · 레나 델타 · 캐나다 북서부(그림과 표에서는 Canada (NW))를 대상으로 하였다(그림~\ref{fig:data}, 표~\ref{tab:data}). 알래스카와 캐나다는 ABoVE 토양 수분 · 활동층 두께 자료 v2\cite{Moore2025}와 CALM 계절 말 격자 값\cite{Streletskiy2025}을, 레나 델타는 ALLena 융해 깊이 자료\cite{Veremeeva2025}를 사용하였다. ABoVE 탐침 측정은 8--10월 값으로 제한하였는데, 그보다 이른 측정은 계절 최댓값과의 차이가 크기 때문이며, 값은 0 cm 초과 300 cm 미만으로 제한하였다. ALLena 는 탐침 한계 표시(probe-limit flag; 탐침이 최대 길이에 도달한 측정)가 없는 8--9월 값으로 제한하고, 위치 · 연도별 계절 최댓값을 라벨로, 그 측정일들의 평균을 관측일로 취하였다. CALM 의 융해 깊이는 융해 계절이 끝나는 8월 말에서 9월 초에 탐침으로 측정하므로\cite{Nelson1997,Brown2000}, 날짜가 없는 CALM 값의 관측일은 연중 244일째(9월 1일)로 두었다. 관측은 좌표를 0.0001$^\circ$(약 10 m)로 반올림한 위치인 지점별 · 연도별로 묶어 평균하고 평균 관측일을 보존하였다. 따라서 ABoVE 측선의 탐침 지점은 각각 별개의 지점이며, 한 1 km 셀 안에 여러 지점이 있을 수 있다. 세 프로그램은 규약(ABoVE 와 ALLena 는 탐침 측선, CALM 은 탐침 약 100개의 1 km 격자 평균)과 지점 선택이 다르며, 모두 보고된 그대로 사용하였다. 라벨, 기후 입력, 예측, 평가는 모두 특정 연도의 한 지점 단위로 정의한다. $y$년의 라벨은 $y$년의 기후와만 짝지어지고, $y$년의 지도는 $y$년의 기후만 쓴다. 연도별 기후 입력이 있는 2010--2024년만 사용하였다. 라벨은 알래스카 \nncyAk{}개 지점-연도(\nyrsAkko), 레나 델타 \nncyLe{}개(\nyrsLeko), 캐나다 \nncyCa{}개(\nyrsCako)이며, 0.5$^\circ$ 블록 \nnblkAk, \nnblkLe, \nnblkCa{}개에 분포한다(표~\ref{tab:data}). 지점은 1 km 셀 \nnkmcellAk{}개, \nnkmcellLe{}개, \nnkmcellCa{}개(셀-연도 \nnkmcyAk{}개, \nnkmcyLe{}개, \nnkmcyCa{}개)에 들어 있고, 한 측선의 지점이 한 셀 · 한 해에 최대 \nnkmmaxAk{}개까지 몰리므로 지점-연도는 1 km 규모에서 독립 표본이 아니며 공간 블록 설계가 이를 다룬다. 자료 출처 · 판 · 약관은 보충 표 S1, 입력 변수는 보충 표 S2, 분석별 사용 연도는 보충 표 S3, 처리 단계는 보충 방법 S1에 정리하였다.

% table1.tex (국문 v15): 표 1. 지역별 라벨과 입력 · 제품의 영역 · 해상도 · 기간 · 사용 연도. 수치는 매크로.
\begin{table}[!htbp]
\caption{자료와 사용 연도. 위: 지역별 ALT 라벨(관측 연도 기후 입력이 있는 2010--2024년 지점-연도; 지점은 좌표를 0.0001$^\circ$ 로 반올림한 관측 위치이고 1 km 셀은 지점을 포함하는 지도 셀이다; 블록은 공간 교차검증과 부트스트랩에 쓴 0.5$^\circ$ 셀이며, 2020년에는 캐나다 라벨이 없다). 아래: 입력과 제품의 공간 영역, 해상도, 제공 기간, 이 연구에서 쓴 연도.}\label{tab:data}
\scriptsize
\setlength{\tabcolsep}{4pt}
\begin{tabular}{@{}lrrrrlrl@{}}
\toprule
지역 & 지점-연도 & 지점 & 1 km 셀 & 블록 & 연도 & 4개 해 이상 지점 & 최다 연도 \\
\midrule
알래스카 & \nncyAk{} & \nncellAk{} & \nnkmcellAk{} & \nnblkAk{} & 2010--2024 & \nnmultiAk{} & \ntopyrAk{} (\ntopyrshareAk\%) \\
레나 델타 & \nncyLe{} & \nncellLe{} & \nnkmcellLe{} & \nnblkLe{} & 2010--2022 & \nnmultiLe{} & \ntopyrLe{} (\ntopyrshareLe\%) \\
캐나다 북서부 & \nncyCa{} & \nncellCa{} & \nnkmcellCa{} & \nnblkCa{} & 2010--2024 & \nnmultiCa{} & \ntopyrCa{} (\ntopyrshareCa\%) \\
\bottomrule
\end{tabular}

\vspace{6pt}
\begin{tabular}{@{}>{\raggedright\arraybackslash}p{0.24\textwidth}>{\raggedright\arraybackslash}p{0.13\textwidth}>{\raggedright\arraybackslash}p{0.12\textwidth}>{\raggedright\arraybackslash}p{0.10\textwidth}>{\raggedright\arraybackslash}p{0.16\textwidth}>{\raggedright\arraybackslash}p{0.15\textwidth}@{}}
\toprule
자료 & 영역 & 해상도 & 제공 기간 & 사용 & 역할 \\
\midrule
ERA5-Land 월별\cite{munozsabater2021_era5_land} & 전 지구 & 0.1$^\circ$(약 9 km) & 1950-- & 2010--2024; 평균 2011--2024 & 기후 입력 \\
Copernicus DEM\cite{CopDEM2021} & 전 지구 & 30 m & 정적 & -- & 지형 \\
SoilGrids 2.0\cite{Poggio2021} & 전 지구 & 250 m & 정적 & -- & 토양 \\
MODIS MOD13Q1, MOD10A2\cite{Didan2021,HallRiggs2021} & 전 지구 & 250 m, 500 m & 2000-- & 2010--2024 & 민감도 시험 \\
ESA CCI 영구동토 확률\cite{Westermann2024} & 북반구 & 1 km & 1997--2021 & 1997--2021 평균 & 지도 영역(확률 $\geq$ 50\%) \\
Wei 등\cite{Wei2026} & 북반구 & 1 km & 연도별 2000--2024 & \nyrsprodWei & 비교 \\
Liu 등\cite{Liu2024} & 북반구 & 1 km & 연도별 2003--2020 & \nyrsprodLiu & 비교 \\
ESA CCI ALT v4\cite{Westermann2024} & 북반구 & 1 km & 연도별 1997--2021 & \nyrsprodCcifour & 비교 \\
ESA CCI ALT v5\cite{Westermann2025} & 북반구 & 1 km & 연도별 1997--2023 & \nyrsprodCcifive & 비교 \\
Aalto 등\cite{Aalto2018} & 환북극 & 30$''$(약 1 km) & 2000--2014 평균 & 셀당 값 하나 & 비교 \\
Yi와 Kimball\cite{YiKimball2019} & 알래스카(ABoVE 영역) & 1 km & 연도별 2001--2015 & \nyrsprodYk & 참고, 미채점 \\
\bottomrule
\end{tabular}
\end{table}


\paperfig{Fig1_data}{1.0}{\textbf{연구 지역, 관측, 자료 기간.} \textbf{a}, ESA CCI 영구동토 확률(1997--2021년 평균)과 세 연구 지역, 2010--2024년 ALT 라벨이 있는 지점 수. \textbf{b}, 관측 연도별 지점-연도 라벨 수. \textbf{c}, 라벨, ERA5-Land 기후 입력, 정적 공변량, 공개 제품 6종의 기간. 짙은 막대는 라벨 지점에 제품 값이 있는 해, 음영은 분석 기간이다.}{fig:data}

\subsection{관측 연도 기후 입력}
2010년 1월--2024년 12월 ERA5-Land 월자료\cite{munozsabater2021_era5_land}(2 m 기온, 표층 토양 온도, 적설 수당량)를 각 지점의 가장 가까운 육지 격자에서 추출하였다. 지점 $c$, 연도 $y$의 융해 도일\cite{Nelson1997,Riseborough2008}은 다음과 같다.
\begin{equation}
\mathrm{TDD}_{c,y} = \sum_{m=1}^{12} \max\!\left(T_{c,y,m},\,0\right) d_{y,m}
\label{eq:tdd}
\end{equation}
여기서 $T_{c,y,m}$은 월평균 기온, $d_{y,m}$은 $m$월의 일수이다. 관측은 대부분 융해 계절이 끝나기 전에 이루어졌으므로, 모델 입력은 관측일 $t$까지 누적한 TDD이다.
\begin{equation}
\mathrm{TDD}^{\mathrm{obs}}_{c,y} = C_{c,y,k-1} + \frac{t - t_{k-1}}{d_{y,k}} \left( C_{c,y,k} - C_{c,y,k-1} \right)
\label{eq:tddobs}
\end{equation}
$C_{c,y,k}$는 $k$월 말까지의 누적 TDD, $k$는 $t$가 속한 달, $t_{k-1}$은 전월의 마지막 날이다. 이와 함께 기후 변수 일곱 개, 즉 연평균 기온, $\sqrt{\mathrm{TDD}}$, 가장 따뜻한 달과 가장 추운 달의 평균 기온, 직전 겨울(전년 10월부터 관측 연도 4월까지)의 동결 도일, 표층 토양 온도 평균, 수문년(전년 10월부터 관측 연도 9월까지) 적설 수당량 평균을 지점 · 연도별로 계산하였다(보충 표 S2, 보충 방법 S2). 자료가 2010년 1월부터 있으므로 2010년의 겨울 변수는 2010년 1--4월(동결 도일)과 1--9월(적설 수당량)로 계산하였다. 같은 지점-연도에서 두 가지 대안 입력, 즉 관측 연도의 계절 전체 TDD 와 연도별 값의 2011--2024년 평균(이하 다년 평균 기후 입력)을 함께 평가하였다. 후자는 기존 제품의 다년 입력에 해당한다. 대부분의 관측이 8월에 이루어지므로, 라벨 지점에서 관측일까지 누적한 TDD 는 계절 전체 합에 대한 비율의 중앙값이 \ntddfracAk, \ntddfracLe, \ntddfracCa(알래스카, 레나 델타, 캐나다)이다. 연도별 지도는 CALM 규약의 계절 말 기준일인 연중 244일째(윤년이 아닌 해의 9월 1일)까지 누적한 TDD 를 쓴다. 이 날짜는 관측일의 합동 사분위 범위(224--234일)보다 10일 늦으나 관측일 전체 범위(213--288일) 안에 있다. 지도 값은 명시된 해의 9월 1일까지 도달한 융해 깊이를 나타낸다.

\subsection{정적 공변량}
지형 변수 여섯 개(고도, 경사, 사면 방향의 사인 · 코사인, 지형 위치 지수, 거칠기)는 30 m Copernicus DEM\cite{CopDEM2021}에서 각 지점 둘레 $33\times33$ 화소 창(약 1 km)으로 계산하였다. 토양 변수 아홉 개(5--15 cm의 점토, 모래, 실트, 용적 밀도, 조립질, pH와 0--5, 5--15, 15--30 cm의 토양 유기탄소)는 250 m SoilGrids 2.0\cite{Poggio2021}에서 추출하였다. 따라서 주 특징 벡터 $\mathbf{x}$는 기후 8개, 지형 6개, 토양 9개의 23개 변수로 이루어진다. 민감도 시험에는 다섯 가지 대안 묶음, 즉 기후만, 기후와 지형, 주 묶음에 2010--2024년 MODIS 여름 NDVI · EVI와 눈 소실 시기(MOD13Q1 · MOD10A2 v6.1)\cite{Didan2021,HallRiggs2021}를 더한 묶음, 주 묶음에 고해상 토양 탄소 · 피복 분류\cite{Palmtag2022}를 더한 묶음, 전년 기후를 포함한 전체 변수를 사용하였다. ESA CCI 영구동토 확률(1997--2021년 평균)\cite{Westermann2024}은 예측 변수가 아니라 지도 범위(확률 50\% 이상)를 정하는 데에만 사용하였다.

\subsection{모델 구성}
Stefan 모델은 융해 깊이를 다음과 같이 나타낸다.
\begin{equation}
\widehat{\mathrm{ALT}}_{\mathrm{S}} = E\,\sqrt{\mathrm{TDD}^{\mathrm{obs}}}
\label{eq:stefan}
\end{equation}
$E$(cm\,($^\circ$C\,d)$^{-1/2}$)는 토양의 열적 성질과 수분을 요약한다\cite{Stefan1891,Riseborough2008}. 이 관계는 레나 델타 관측 지점에서 계절 안에 잘 성립하며, 반복 관측된 셀에서는 연도 간에도 더 큰 산포와 함께 성립한다(보충 그림 S1). Stefan 모델을 물리 기준값으로 택한 이유는 매개변수가 하나여서 소수의 관측으로 다시 맞춰 지역 간에 전이할 수 있고 입력이 연도별 TDD 이기 때문이다. 척도 계수를 둔 Kudryavtsev 모델\cite{Kudryavtsev1974,RomanovskyOsterkamp1997}과 두 모델의 비음수 결합은 대안 기준값으로 평가하였다(보충 표 S4). 흔히 edaphic factor 라 부르는 이 계수를 지역 관측으로 맞추는 것은 Stefan 기반 ALT 지도화의 관행이다\cite{Nelson1997,ShiklomanovNelson2002}. $n$개 학습 라벨 $y_i$와 $s_i = \sqrt{\mathrm{TDD}^{\mathrm{obs}}_i}$에 대해 계수는 원점을 지나는 최소제곱으로 재보정한다.
\begin{equation}
\hat{E} = \frac{\sum_{i=1}^{n} y_i\, s_i}{\sum_{i=1}^{n} s_i^{2}}
\label{eq:E}
\end{equation}
원천 계수 Stefan 모델은 다른 두 지역에서 맞춘 $\hat{E}$를, 재보정 Stefan 모델은 대상 지역 학습 라벨로 맞춘 $\hat{E}$를 쓴다. 물리 잔차 모델은 재보정 Stefan 예측을 기준값으로 두고 학습한 보정량을 더하며, 이는 물리 기반 학습의 잔차 모델링 형태이다\cite{Willard2022}(그림~\ref{fig:workflow}).
\begin{equation}
\widehat{\mathrm{ALT}} = \hat{E}\, s + \lambda\, g(\mathbf{x}), \qquad g = \arg\min_{f} \sum_{i=1}^{n} \left[ \left( y_i - \hat{E}\, s_i \right) - f(\mathbf{x}_i) \right]^{2}
\label{eq:residual}
\end{equation}
$g$는 CatBoost gradient boosting 회귀기\cite{Prokhorenkova2018}(반복 600, 학습률 0.03, 깊이 6, L2 정칙화 3)이며, $\lambda=0.5$는 보정량을 물리 기준값 쪽으로 수축한다. CatBoost 를 기본 학습기로 쓴 이유는 gradient boosting 트리가 특징 척도 조정이 필요 없고 이 규모의 표 자료에서 통상의 기준이기 때문이다\cite{Grinsztajn2022}. 랜덤 포레스트, 다층 퍼셉트론, TabM, FT-Transformer\cite{Gorishniy2021,Gorishniy2025,Holzmuller2024}로 바꾸어 민감도를 시험한 결과, 어느 학습기도 세 지역 모두에서 가장 정확하지는 않았다(보충 그림 S2). $\lambda=0.5$ 는 지역 안 이득(공간 블록 검증에서 재보정 Stefan 모델 대비 RMSE 감소량; $\lambda=1$ 에서 레나 델타 \ndlamOneLeblk{} cm)과 지역 홀드아웃(한 지역 전체를 학습에서 제외하는 검증, 아래 정의) 오차($\lambda=1$ 에서 레나 델타 \ndlamOneLereg{} cm)가 모두 $\lambda$ 와 함께 커지므로 둘 사이의 절충으로 골랐고, $\lambda$ = 0.25 와 1 은 민감도로 보고한다. 비교한 세 구조 가운데 직접 ML 은 물리 정보를 쓰지 않고, 물리 입력 ML 은 물리 예측을 입력 특징으로, 물리 유사라벨 모델은 물리 예측을 추가 학습 목표로 쓴다.
\begin{equation}
\begin{aligned}
\widehat{\mathrm{ALT}}_{\mathrm{D}} &= f_{\mathrm{D}}(\mathbf{x}), \qquad
\widehat{\mathrm{ALT}}_{\mathrm{I}} = f_{\mathrm{I}}\!\left(\mathbf{x},\, \hat{E}\,s,\, \widehat{\mathrm{ALT}}_{\mathrm{K}}\right), \\
\widehat{\mathrm{ALT}}_{\mathrm{A}} &= f_{\mathrm{A}}(\mathbf{x}) \;\text{학습 자료}\; \{(\mathbf{x}_i, y_i)\}_{i=1}^{n} \cup \{(\mathbf{x}_j, \hat{E}\,s_j)\}_{j=1}^{n}
\end{aligned}
\label{eq:alt}
\end{equation}
직접 ML($f_{\mathrm{D}}$)은 $\mathbf{x}$만으로 ALT를 학습하고, 물리 입력 ML($f_{\mathrm{I}}$)은 재보정 Stefan 예측과 Kudryavtsev\cite{Kudryavtsev1974} 예측 $\widehat{\mathrm{ALT}}_{\mathrm{K}}$를 특징으로 더하며, 물리 유사라벨 모델($f_{\mathrm{A}}$)은 학습 자료 $n$개에, 평가 겹에 속한 셀에서 뽑은 같은 수의 표본 $j$를 더하여 학습한다. 더한 표본의 입력은 그 셀의 공변량이고, 목표값은 학습 겹 자료로 맞춘 Stefan 예측이다. 모든 학습 모델은 난수 시드 3개로 학습하여 예측을 평균하였다(대상 라벨이 없는 전이 모델은 시드 1개). 지역별 전체 라벨로 학습한 물리 잔차 모델과 직접 ML 에 대해, 각 입력 변수가 예측에 기여한 정도인 SHAP(SHapley Additive exPlanations) 값\cite{Lundberg2017}을 계산하였다.

\paperfig{Fig2_workflow}{1.0}{\textbf{지점-연도 설계, Stefan 기준값, 물리 잔차 모델, 검증 설계.} \textbf{a}, 관측일의 융해 깊이와 계절 말의 ALT(왼쪽), 레나 델타 셀의 2013년 누적 융해 도일과 관측일까지의 모델 입력 · 9월 1일까지의 지도 입력(가운데), 알래스카 지점-연도의 관측 ALT 대 $\sqrt{\mathrm{TDD_{obs}}}$와 Stefan 기준값 · 잔차 5개(오른쪽). \textbf{b}, 물리 잔차 모델: Stefan 기준값에, 기후 · 지형 · 토양 입력 23개로 잔차를 학습한 gradient boosting 보정량을 더한다(식~(\ref{eq:residual})). 알래스카 2014년 지도로 예시하였다. \textbf{c}, 레나 델타 라벨 셀의 검증 설계(색은 겹, 관측 연도 또는 시험 지역)와 알래스카 지도의 2012년 · 2019년 ALT 편차.}{fig:workflow}

\subsection{검증 설계와 불확실성}
모든 모델은 다섯 가지 검증 설계에서 겹 밖 예측(out-of-fold prediction; 각 지점-연도를 그것이 속한 시험 겹을 제외한 자료로 학습한 모델로 예측한 값)으로 평가하였다\cite{Roberts2017,Ploton2020}. 다섯 설계는 지점-연도 단위 무작위 5겹 교차검증, 겹마다 지점-연도 수가 비슷하도록 0.5$^\circ$ 블록을 배정한 공간 블록 5겹 교차검증(주 설계), 지점-연도 20개 이상인 각 관측 연도를 나머지 연도로 학습한 모델로 예측하는 연도 홀드아웃, 평가 대상 지점-연도의 블록 겹과 연도를 모두 학습에서 제외하는 블록-연도 홀드아웃(지점-연도 10개 이상인 겹), 시험 지역에서 100 km 안의 학습 셀을 뺀 뒤 다른 두 지역으로 학습한 모델로 예측하는 지역 홀드아웃이다. 정확도는 다음과 같이 나타낸다.
\begin{equation}
\mathrm{RMSE} = \left[ \frac{1}{n} \sum_{i=1}^{n} \left( \hat{y}_i - y_i \right)^{2} \right]^{1/2}, \qquad \Delta\mathrm{RMSE}_{A,B} = \mathrm{RMSE}_{A} - \mathrm{RMSE}_{B}
\label{eq:rmse}
\end{equation}
$\Delta\mathrm{RMSE}$의 두 항은 같은 지점-연도에서 계산하며, 음수이면 방법 $A$ 의 오차가 작다. 불확실성은 짝지은 블록 부트스트랩\cite{Efron1993,Roberts2017}으로 정량화하였다. 0.5$^\circ$ 블록을 복원 추출로 $B=2000$회(라벨 수 실험과 공변량 묶음 · 학습기 민감도 실행의 일부는 1000회) 재추출하고, 재추출된 블록의 모든 지점-연도로 $\Delta\mathrm{RMSE}$를 다시 계산하여, 반복값의 2.5와 97.5 백분위수를 95\% 구간으로 정하였다.
\begin{equation}
\mathrm{CI}_{95} = \left[ Q_{0.025}\!\left( \Delta\mathrm{RMSE}^{(b)} \right),\; Q_{0.975}\!\left( \Delta\mathrm{RMSE}^{(b)} \right) \right], \quad b = 1,\ldots,B
\label{eq:ci}
\end{equation}
구간이 0을 제외하면 차이가 판정된 것으로 기술한다. 한 1 km 셀에 여러 지점이 들어갈 수 있으므로(셀당 지점 수 중앙값: 알래스카 \ncySitesMedAk, 레나 델타 \ncySitesMedLe, 캐나다 \ncySitesMedCa{}개), 주요 비교는 관측과 예측의 셀-연도 평균으로도 채점하였다(보충 표~S5). 이득이 나타나는 조건은 공변량 여덟 개(관측 연도의 계절 전체 TDD, 동결 도일, 적설 수당량, 토양 유기탄소, 점토, 경사, 고도, 여름 NDVI) 각각을 지역 안에서 세 등분한 구간(이하 삼분위 구간)별로 분석하였고, 지점-연도 30개 또는 블록 5개 미만인 구간은 평가하지 않았다. 연도별 RMSE 는 지점-연도 30개 이상인 해에서 계산하였다(보충 그림 S7).

\subsection{라벨 수 · 전이 실험}
지역 안 실험에서는 공간 블록 겹마다 학습 라벨을 $n$ = 25, 50, 100, 250, 500, 1000, 2500개 또는 전량으로 줄이고(무작위 3회), 같은 $n$개 라벨로 Stefan 계수를 재보정한 뒤 시험 겹 전체를 채점하였다. 전이 실험에서는 다른 두 지역의 전체 라벨과 대상 지역 학습 겹에서 뽑은 $n$ = 0, 25, 50, 100, 250, 500개 라벨로 학습한 모델로 대상 지역의 공간 블록 겹을 채점하였다. Stefan 계수는 대상 라벨 $n$개로 다시 맞추었고, 잔차 학습기는 원천 계수 기준값으로 잔차를 계산한 원천 지역 자료와 다시 맞춘 기준값으로 잔차를 계산한 대상 지역 자료를 함께 써서 학습하였다.
\begin{equation}
\begin{aligned}
\widehat{\mathrm{ALT}}_{\mathrm{T}} &= \hat{E}_{\mathrm{T}}\, s + \lambda\, g_{\mathrm{ST}}(\mathbf{x}), \\
g_{\mathrm{ST}} &= \arg\min_{f} \Big[ \sum_{i \in \mathrm{S}} \big( y_i - \hat{E}_{\mathrm{S}}\, s_i - f(\mathbf{x}_i) \big)^{2} + \sum_{i \in \mathrm{T}} \big( y_i - \hat{E}_{\mathrm{T}}\, s_i - f(\mathbf{x}_i) \big)^{2} \Big]
\end{aligned}
\label{eq:transfer}
\end{equation}
S 는 원천 지역 자료, T 는 대상 지역 자료 $n$개, $\hat{E}_{\mathrm{S}}$는 원천 자료로 맞춘 계수, $\hat{E}_{\mathrm{T}}$는 대상 자료로 다시 맞춘 계수이다($n < 3$이면 $\hat{E}_{\mathrm{T}} = \hat{E}_{\mathrm{S}}$). 이 실험에서는 대상 지역에서 100 km 안의 원천 지역 자료를 제외하지 않았다.

\subsection{공개 제품 비교}
공개 저장소에 있는 제품 6종을 내려받아 라벨 지점에서 추출하였다(최근접 1 km 화소; 보충 방법 S3, 표 S1). 북반구 제품은 Wei 등\cite{Wei2026}(1 km, 연도별 2000--2024년), Liu 등\cite{Liu2024}(1 km, 연도별 2003--2020년), ESA CCI Permafrost ALT v4(1997--2021년)와 v5(1997--2023년)\cite{Westermann2024,Westermann2025}, Aalto 등\cite{Aalto2018}(2000--2014년 기준 기간, 셀당 값 하나)이고, 알래스카 전용 제품인 Yi와 Kimball\cite{YiKimball2019}(2001--2015년)은 한 지역만 덮고 2015년에 끝나므로 참고로만 쓰고 채점하지 않았다. 연도별 제품은 관측 연도의 값을 사용하여 제품과 모델을 같은 지점 · 같은 해에서 비교하였다. 제품과 모델은 Wei 등, Liu 등, ESA CCI v4가 모두 관측 연도 값을 가진 공통 지점-연도(\nyrscc{}년)에서 비교하였다. ESA CCI v5와 Aalto 등은 값이 있는 공통 지점-연도에서 채점하였고, 연도별 비교에는 Wei 등 값이 있는 그 해의 모든 지점-연도를 썼다. Wei 등은 학습 지점 좌표를 공개하였으므로, 학습 자료와 겹치는 영향을 줄이기 위해 이 지점에서 25 km 안의 라벨 셀을 제외한 비교(이하 누설 마스크)를 함께 수행하였다. Liu 등은 CALM 과, 이 연구에서 쓴 알래스카 · 캐나다 관측의 이전 판인 ABoVE 토양 수분 · 활동층 자료로 랜덤 포레스트를 학습하였으므로\cite{Liu2024} 그 지점에서의 제품 값은 우리 라벨과 독립이 아니며, Wei 등의 학습 지점도 같은 방식으로 우리 알래스카 지점과 겹친다. 잔차 학습기가 제품도 보정할 수 있는지 시험하기 위해, 관측 연도의 Wei 등 제품 값 $p$를 학습 겹에서 다음과 같이 재보정하고
\begin{equation}
\hat{p} = a + b\, p
\label{eq:prod}
\end{equation}
잔차 $y-\hat{p}$를 식~(\ref{eq:residual})과 같은 공변량, 학습기, $\lambda$, 겹으로 학습하였다(보충 방법 S4).

\subsection{연도별 지도, 적용 가능 영역, 예측 구간}
세 지역 각각에서 재보정 Stefan 모델과 물리 잔차 모델을 그 지역의 모든 지점-연도로 학습하였다. 두 모델은 ESA CCI 영구동토 확률 50\% 이상인 1 km 셀, 즉 Obu 등\cite{Obu2019}의 분류에서 활동층이 영구동토 위에 존재하고 라벨이 분포하는 연속 · 불연속 영구동토 구역에 속하는 셀 가운데 수면(Global Surface Water 평균 출현율 50\% 이상)과 입력 결측 셀을 뺀 셀(\nncellmapAk{}개, \nncellmapLe{}개, \nncellmapCa{}개)에 2011--2024년 각 연도의 기후(9월 1일까지 누적한 TDD)를 넣어 적용하였다. 적용 가능 영역은 Meyer와 Pebesma\cite{MeyerPebesma2021}를 따랐다. $\mathbf{z}$를 표준화한 23차원 특징 벡터(2011--2024년 평균 기후)라 할 때, 지도 셀 $k$의 비유사도는 모든 학습 행 $j$에 대한 $d_k = \min_{j} \lVert \mathbf{z}_k - \mathbf{z}_j \rVert_2$이며, 다음을 만족하면 적용 가능 영역 안으로 본다.
\begin{equation}
d_k \leq \tau, \qquad \tau = Q_{0.95}\!\left\{ \min_{j:\, b(j) \neq b(i)} \lVert \mathbf{z}_i - \mathbf{z}_j \rVert_2 \;:\; i \in \text{학습 행} \right\}
\label{eq:aoa}
\end{equation}
$b(\cdot)$는 공간 블록이다. 연도 $y$, 셀 $k$의 편차는 $a_{k,y} = \widehat{\mathrm{ALT}}_{k,y} - \overline{\widehat{\mathrm{ALT}}}_{k}$이며(윗줄은 2011--2024년 평균), 지역 평균 편차는 지도 셀 전체의 $a_{k,y}$ 평균이고, 셀별 추세는 최소제곱 기울기이다.
\begin{equation}
\beta_k = \frac{\sum_{y} \left( y - \bar{y} \right) a_{k,y}}{\sum_{y} \left( y - \bar{y} \right)^{2}}
\label{eq:trend}
\end{equation}
단위는 10년당 cm이며, 이는 이전 연도별 ALT 제품의 방식과 같다\cite{Liu2024,YiKimball2019}. 지도 값의 90\% 예측 구간은 다음과 같다.
\begin{equation}
\left[ \widehat{\mathrm{ALT}}_{k,y} + Q_{0.05}(r),\; \widehat{\mathrm{ALT}}_{k,y} + Q_{0.95}(r) \right]
\label{eq:pi}
\end{equation}
$r$은 그 지역의 공간 블록 겹 밖 잔차이다. 이 구간은 분할 등각 예측(split-conformal prediction) 방식의 구간\cite{Lei2018}으로 지역 안에서 폭이 일정하며, 셀별 불확실성은 식~(\ref{eq:aoa})의 비유사도 $d_k$ 로 나타낸다. 표시용으로 1 km 값을 250 m(레나 델타) 또는 500 m(알래스카, 캐나다) 음영기복도 격자에 쌍선형 보간하였다. 지역별로 우리 지도와 두 제품 모두 값이 있는 관측 1 km 셀이 가장 많은 해(알래스카 · 캐나다 2013년, 레나 델타 2016년)의 지도를 두 ML 제품과 비교하였다. 그 셀에서 같은 해의 제품 값, 관측의 셀 평균, 우리 모델의 공간 블록 겹 밖 예측을 썼고, 제품 지도는 우리 지도의 셀에만 표시하며, 지역마다 색 범위는 적용 가능 영역 안 이 연구 지도 값의 5--95 백분위수를 제품 중앙값까지 확장하여 정하고 범위를 넘는 제품 값은 범위 끝 색으로 표시한다.

\subsection{계산}
모델은 난수 시드를 고정해 96코어 CPU 노드에서, 신경망 학습기는 NVIDIA A100 GPU 1대에서 학습하였고, 자료 준비 · 집계 · 그림은 로컬 서버에서 수행하였다(보충 방법 S5).

\subsection{코드 가용성}
자료 처리, 모델 적합, 평가에 쓴 코드는 [저장소 · DOI]에서 받을 수 있다. 본문의 모든 수치는 저장된 결과 표에서 스크립트로 자동 생성하였다.
